\cotefile{91}
\scanpage{59}
\unsure{comme par|comme un|canonique}
\[
T\longrightarrow f_*^{\mathrm{ét.}}(T)
\]
qui est tel que l'hom correspondant $\underline{\mathrm{Hom}}_S(\overbrace{f_*^{\mathrm{ét}}(T)}^{\bar T},S')\longrightarrow\underline{\mathrm{Hom}}_S(T,S')$ est un isom. pour tt $S'$. (\unsure{Préciser|Précisions}\,\dots). Il en résulte que pour tt \deleted{$s\in S,$} \unsure{changemt|changement} de base (p. \unsure{ex.} \deleted{$\underline{\mathrm{Hom}}_{k(s)}(\bar T_s,S'_s)\longrightarrow$} \unsure{de} $S=\operatorname{Spec}(\overline{k(s)})$] la $\unsure{=|\simeq}$ relation reste vraie\,; faisant $S'=S\times\underline I$ ($I$ un fini) on en conclut que $T\to f_*^{\mathrm{ét}}(T)$ est surjectif et \inserted{(si $T/S$ est fini)} \unsure{géométriquement} injectif (regarder \deleted{\unsure{sur}} \inserted{au dessus de} chaque pt $s\in S$).

Inversement, si $T\to\bar T$ est un morphisme surjectif et géom injectif de $T$ (fini \inserted{et plat} sur $S$) sur $\bar T$ rev. étale de $S$, alors on a $\underline{\mathrm{Hom}}_S(\bar T,S')=\mathrm{Hom}_S(T,S')$ pour tout $S'$ non ramifié sur $S$ [\unsure{a fortiori|et fini|et finiment}] et comme la $=$ chose reste vrai par ext. de la base \inserted{$S$}, on conclut \deleted{\unsure{Ho}} $\underline{\mathrm{Hom}}_S(\bar T,S')\xrightarrow{\sim}\underline{\mathrm{Hom}}_S(T,S')$, i.e. \unsure{Term.|Terminé|$T$ \illegible{1}}

% archive sheet number: 58
\scanpage{61}
\textcircled{1}
\uline{Set up général}

\deleted{\unsure{Soi}}Soit $S$ un préschéma. \deleted{\unsure{Définition des foncteurs}}

Soit $\unsure{\underline{\underline{\mathbb W}}|\underline{\mathbb W}|\underline W}$ un foncteur covariant de la catégorie des $S$-préschémas dans la catégorie des \deleted{$\unsure{S|E}$} préschémas. On suppose $\underline{\underline{\mathbb W}}$ ``de nature locale'', i.e.

a i) Si $i:T\longrightarrow T'$ est une immersion ouverte, alors $\underline{\underline{\mathbb W}}(i):\underline{\underline{\mathbb W}}(T)\to\underline{\underline{\mathbb W}}(T')$ est une immersion ouverte. [Donc si $U$ ouvert dans $T$, on identifie $\underline{\underline{\mathbb W}}(U)$ à un ouvert de $\underline{\underline{\mathbb W}}(T)$]

% slanted
\margin{\unsure{on|ou} \unsure{peut faire|pour faire} cette hypothèse \unsure{là|?} \deleted{\unsure{tout}} de suite}

a ii) $\underline{\underline{\mathbb W}}(\bigcup U_i)=\bigcup_i\underline{\underline{\mathbb W}}(U_i)$ pour une réunion \unsure{quelc.} de parties ouvertes de $T$

\[
\underline{\underline{\mathbb W}}(U\cap V)=\underline{\underline{\mathbb W}}(U)\cap\underline{\underline{\mathbb W}}(V)\tag*{a iii)}
\]

\uline{Exemple} I. Soit $S'$ un $S$-préschéma. Posons $\underline{\underline{\mathbb W}}_{S'}(T)=T\times_S S'$ : on obtient ainsi \inserted{$\underline{\underline{\mathbb W}}_{S'}$} un foncteur covariant de nature locale des $S$-préschémas dans les $S'$-préschémas (à fortiori dans les préschémas).

Cas I' : $S'$ est fini et loc. libre sur $S$

\uline{Exemple} II Soit $W$ un \deleted{\unsure{$S$|$A$}} préschéma en anneaux commutatifs sur $S$, supposons que pour tt $T/S$, le faisceau d'anneaux $\underline\Gamma(W_T/T)$ sur $T$ fasse de $T$ un préschéma (cf exemple plus bas). \deleted{On a ainsi}
% slanted
\margin{annelé en anneaux locaux}
% préschéma circled
% archive sheet number: 60
\scanpage{62}
\deleted{\unsure{un fonc}} et que l'hom naturel $\underline\Gamma(W_T/T)\unsure{-|\to}\underline\Gamma(W_{T'}/T')$ pour $T'\to T$ \deleted{\unsure{de fa|défini}} soit \uline{local} sur les fibres en faisceaux. On obtient ainsi un foncteur \inserted{$\unsure{\underline{\mathbb W}|\underline{\underline{\mathbb W}}}$} de nature \unsure{locale|covt|covariant} de $S$-préschémas en $S$-préschémas, et un isom fonctoriel $|\underline{\mathbb W}(T)|\simeq|T|$ (espaces ss-jacents). \deleted{\unsure{on peut évidemment|on voit évidemment|vraisemblablement} d'ailleurs $\Lambda$, puisque \unsure{l'on aurait} $\mathrm{Hom}_S(T,\Lambda)=$} \inserted{Nous donnerons des} conditions suffisantes plus \unsure{loin|bas} pour que $\Lambda$ ait les \unsure{propr. voulues|voisinage}.

\uline{Exemple} III Soit $S'=\operatorname{Spec}(A)$, où $A$ est un faisceau loc. libre d'algèbres comm. sur $\underline O_S$. Supposons de plus que $S'\unsure{-|\to}S$ soit un homéomorphisme universel. Posons $\Lambda=\deleted{\unsure{\operatorname{Spec}}}\mathbf V(\check A)$. On est alors simultanément dans le cas de l'exemple I et de l'exemple II.
\margin{$S'\to S$}

\uline{Exemple} III' Comme II, mais on suppose donnée une augmentation $W\longrightarrow S[t]$ d'anneaux, telle que pour tout $T/S$, l'augmentation correspondante $\underline O_{\underline{\mathbb W}(T)}\longrightarrow\underline O_T$ soit surjective, de noyau nilpotent.
% slanted
\margin{\unsure{blague sur} exemple II : foncteurs de Greenberg}

Soit $\unsure{\mathfrak X|\mathcal X|\underline X}$ un préschéma sur $\underline{\mathbb W}(S)$. Considérons le foncteur contravariant en le $S$-préschéma $T$
\[
\mathrm{Hom}_{\mathbb W(S)}(\mathbb W(T),\mathfrak X)\tag*{(*)}
\]
% archive sheet number: 61
\scanpage{63}
On \deleted{\unsure{voit|veut}} se demande s'il est représentable, i.e. on demande la valeur en $T$ du \uline{foncteur adjoint} (au sens de Kan) du foncteur $\underline{\mathbb W}$ des $S$-préschémas en $T$-préschémas. Une solution du Pb universel posé par le foncteur en $T$ ci-dessus est par définition un couple
\[
 X\quad(S\text{-préschéma})\qquad \underline{\mathbb W}(X)\xrightarrow{\alpha}\unsure{\mathfrak X|\mathcal X}
\]
tel que pour tout $T$, l'application naturelle
\[
 \mathrm{Hom}_S(T,X)\xrightarrow{\sim}\mathrm{Hom}_{\mathbb W(S)}(\mathbb W(T),\mathfrak X)
\]
\deleted{soit une} \unsure{définie|obtenue} par factorisation
\[
 \underline{\mathbb W}(T)\xrightarrow{\mathbb W(u)}\underline{\mathbb W}(X)\xrightarrow{\alpha}\mathfrak X
\]
soit \uline{bijective}. Comme le foncteur $(*)$ en $T$ est de ``nature locale'', il suffit d'ailleurs de vérifier la bijectivité pour $T$ affine \deleted{\unsure{et}} \inserted{sur un ouvert} affine \deleted{\unsure{de|sur}} $S$, p. ex.

Désignons par $\underline{\mathbb W}^a$ le foncteur adjoint de $\underline{\mathbb W}$ (défini sur une ss-catégorie de la catégorie des $T$-préschémas)
\[
 \mathrm{Hom}_{S^a}(T^a,\mathfrak X)\simeq\mathrm{Hom}_S(T,\mathfrak X
\]
% archive sheet number: 62
\scanpage{64}
(i) Si $\deleted{\unsure{\mathcal O}}\unsure{\underline{\mathbb W}|\underline W}^a(\unsure{\mathfrak X|\mathcal X})$ et $\underline{\mathbb W}^a(\unsure{\mathfrak Y|\mathcal Y})$ sont définis, il en est de $\unsure{=|\simeq}$ de $\mathbb W^a(\mathfrak X\times_{\mathbb W(S)}\mathfrak Y)$, et on a
\[
 \mathbb W^a(\mathfrak X\times_{\mathbb W(S)}\mathfrak Y)\simeq \mathbb W^a(\mathfrak X)\times_S \mathbb W^a(\mathfrak Y)
\]
et on \deleted{\unsure{a}} plus généralement \deleted{\unsure{fon|pour}} la formule analogue pour les limites projectives finies [\unsure{donc|d'où} $(\alpha)$ l'objet unité $(\beta)$ les produits fibrés, d'où comme conséquence $(\gamma)$ Noyaux]

\[
 W^a(\mathfrak X)\times_S T\simeq W_T^a(\mathfrak X\times_{\mathbb W(S)}\mathbb W(T))\tag*{(ii)}
\]
[l'existence du premier \unsure{mb|m\textsuperscript{b}} impliquant celle du second\,; la définition de $W_T$ est claire \unsure{\dots}]

De (i) résulte que le foncteur $\mathbb W^a$ transforme un préschéma en groupes, anneaux etc en préschéma du $=$ type.
% archive sheet number: 63
\scanpage{65}
\uline{Supposons} que $\mathbb W^a(\mathbb W(S)[t])$ existe. C'est alors de façon naturelle un préschéma en anneaux sur $S$, que nous noterons $W$, appelé \uline{préschéma d'anneaux générateur du foncteur} $\mathbb W$.

[Dans l'exemple II, c'est le schéma d'anneaux $W$ de départ. Dans le cas de l'exemple I, \deleted{c'est le} si \margin{\unsure{i)|')}} \deleted{schéma} on suppose $S'$ \inserted{fini et} loc. libre sur $S$, c'est le \deleted{\unsure{faisc|faisc.}} schéma en anneaux $V(\check{\underline A})$, où $\underline A$ est le faisceau d'algèbres sur $S$ associé à $S'$].

\margin{\unsure{')|i)}}Soit alors $\unsure{\mathfrak X|\mathcal X}$ ss-schéma fermé de $\mathbb W(S)[t_1,\ldots,t_r]$ défini par \deleted{\unsure{un}} le faisceau d'idéaux dans \margin{\unsure{$\alpha$|a}} $\underline O_{\mathbb W(S)}[t_1,\ldots t_r]$ engendré par des sections $F_i$ $(i=1,\ldots s)$, donc $\mathfrak X$ est le produit fibré, sur $\unsure{\mathfrak Y|\mathcal Y}^s$, de $\mathfrak Y^r$ et $\mathfrak Y^0$ [où $\unsure{\underline{\mathfrak Y}|\mathfrak Y|\mathcal Y}\deleted{^{\unsure{*|1}}}=\mathbb W[S][t]$).
% archive sheet number: 64
\scanpage{66}
Alors $\mathbb W^a(\unsure{\mathfrak X|\mathcal X})$ existe, et est ainsi donné comme produit fibré. De façon précise, les $F_j$ \inserted{$(1\leq j\leq s)$} définissent des polynômes dans les $t_i$ $(1\leq i\leq r)$, à coefficients dans \deleted{\unsure{les}} $\Gamma(\Lambda/S)$, d'où \deleted{\unsure{des mor}} un $S$-morphisme $\Lambda^r\longrightarrow\Lambda^s$, et $\mathbb W^a(\mathfrak X)$ est $\simeq$ l'image inverse de la section nulle de $\Lambda^s$.

\uline{Exemples} \deleted{$\unsure{\mathfrak Y}$}, $(G_{a_{\underline{\mathbb W}(S)}})^n$, $\mathrm{Gl}(\unsure{m|n})_{\mathbb W(S)}$, etc.

Soit $\underline{\mathbb W}'(T)\xrightarrow{\deleted{\unsure{\varphi}}\alpha(T)}\underline{\mathbb W}^{\deleted{\unsure{'|*}}}(T)$ un morphisme fonctoriel de préschémas, définissant en particulier
\[
 \underline{\mathbb W}'(S)\unsure{-|\to}\underline{\mathbb W}(S)
\]
Soit de plus $\mathfrak X$ sur $\mathbb W(S)$, posons
\[
 \mathfrak X'=\mathfrak X\times_{\underline{\mathbb W}(S)}\underline{\mathbb W}'(S)
\]
On a alors, pour tout $T$, un hom \deleted{\unsure{cf. identité}}
\[
 \mathrm{Hom}_{\underline{\mathbb W}(S)}(\underline{\mathbb W}(T),\mathfrak X)\longrightarrow\mathrm{Hom}_{\underline{\mathbb W}'(S)}(\underline{\mathbb W}(T)\times_{\underline{\mathbb W}(S)}\underline{\mathbb W}'(S)\deleted{\unsure{)}},\mathfrak X')\tag*{(*)}
\]
(image réciproque d'un morphisme par $\mathbb W'(S)-\mathbb W(S)$]

d'autre part on a
\[
 \mathbb W'(T)\to\mathbb W(T)\times_{\mathbb W(S)}\mathbb W'(S)
\]
\deleted{\unsure{par fonct|par produit}} grâce aux composantes $\underline{\mathbb W}'(T)-\deleted{\mathbb W}\underline{\mathbb W}(T)$ et $\mathbb W'(T)\to\mathbb W'(S)$, d'où
\[
 \mathrm{Hom}_{\mathbb W'(S)}(\mathbb W(T)\times_{\mathbb W(S)}\mathbb W'(S),\mathfrak X')\to\mathrm{Hom}_{\mathbb W'(S)}(\mathbb W'(T),\mathfrak X')\tag*{(**)}
\]
d'où en composant (*) et (**) un \deleted{\unsure{hom}} \inserted{morphisme}, évidemment fonctoriel en $T$
\[
 \mathrm{Hom}_{\mathbb W(S)}(\mathbb W(T),\mathfrak X)\longrightarrow\mathrm{Hom}_{\mathbb W'(S)}(\mathbb W'(T),\mathfrak X')
\]
Par suite, si $\mathbb W^a(\mathfrak X)$ et $W'{}^a(\mathfrak X')$ existent on a un morphisme canonique
\[
 \mathbb W^a(\mathfrak X)\xrightarrow{\alpha^*(\mathfrak X)}W'{}^a(\mathfrak X\deleted{^{\unsure{\prime|)}}}\times_{\mathbb W(S)}\mathbb W'(S))
\]
% archive sheet number: 65
\scanpage{67}
De cette façon, pour $\underline{\mathbb W}$ fixé, \deleted{$\unsure{\mathfrak X|\mathcal X}$ fixé} \deleted{$\mathbb W'$ \unsure{var}} $\mathfrak X$ et $\mathbb W'$ variables,
\[
 \mathbb W'{}^a(\mathfrak X\times_{\mathbb W(S)}\mathbb W'(S))
\]
devient un \uline{bifoncteur} en $\mathfrak X,\mathbb W'$, \uline{covariant} en $\mathfrak X$, \uline{contravariant en} $\mathbb W'$

\textcircled{2} \uline{Critère d'existence} : \uline{Cas affine}

% slanted
\margin{Ici, il suffit que $\mathbb W$ \unsure{au sens|au dessus} transforme affines en affines \deleted{\unsure{et soit fonc}}}
% slanted
\margin{Foncteur $\mathbb W$ \uline{spécial} [\unsure{associé} à un schéma \deleted{\unsure{d'}} \inserted{en} anneaux sur $S$ affine sur $S$]}

Supposons que $\underline{\mathbb W}$ transforme schémas affines (sur ouverts affines de $S$) en schémas affines. Donc, au dessus d'un ouvert affine de $S$, d'anneau $A$, $\underline{\mathbb W}$ est défini par un foncteur covariant, également noté $\underline{\mathbb W}$, transformant $A$-algèbres \inserted{$B$} en anneaux $\underline{\mathbb W}(B)$. [\unsure{Avec|Pour} l'hypothèse de \deleted{localité} \inserted{la nature locale} de ce foncteur, que nous n'\unsure{expliciterons|explicitons|explicitation} pas]. On suppose qu'on a $\underline{\mathbb W}(B)=\mathrm{Hom}_{A\text{-alg.}}(\Lambda,B)$ où $\Lambda$ est \deleted{l'algèbre affine} \inserted{l'anneau} d'un schéma en anneaux \deleted{affine} sur $S$, affine sur $S$.
% archive sheet number: 66
\scanpage{68}
[Cela est le cas en particulier dans le cas de l'exemple I, \deleted{$S'$ étant} \inserted{lorsque $S'$ est} fini et loc. libre sur $S$, et l'exemple II, \inserted{lorsque} $W$ \deleted{\unsure{de|est}} est affine sur $S$].

\uline{Proposition} Sous ces conditions, si \deleted{$\unsure{\mathfrak X|\mathcal X}$ est} $\mathfrak X$ est un schéma \uline{affine} sur $\underline{\mathbb W}(S)$, $\mathbb W^a(\mathfrak X)$ existe et est un schéma \uline{affine} sur $S$.

\uline{Démonstration} \deleted{Soit $\mathfrak X$} \inserted{On peut supposer $S$ affine.} Si $\mathfrak X=\operatorname{Spec}(\unsure{\mathcal A|\mathfrak A|\mathcal B|\mathcal C})$ où $\mathcal A$ est une algèbre de type fini \deleted{\unsure{de}} sur $\underline{\mathbb W}(A)$, engendrée par $r$ éléments, il résulte des \unsure{sorites|résultats} du \S{} N\textsuperscript{o} 1 que $\mathbb W^a(\mathfrak X)$ existe et est $\simeq$ un sous-préschéma fermé de $\mathbb W^r$ [défini par des ``équations \unsure{polynômiales|polynôminales|polynomiales} \unsure{à|:}\inserted{\unsure{de|dans}} coeff. dans $\Gamma(\mathbb W/S)$]. Comme $\mathbb W$ est affine, il en est de $\unsure{=|\simeq}$ de \deleted{$\mathbb W$} $\mathbb W^r$ donc aussi de $\mathbb W^a(\mathfrak X)$ (\unsure{ouf|cqf|vif|cqfd}). On a obtenu d'ailleurs

\uline{Cor}. \deleted{Si} \deleted{\unsure{Alors} $\mathcal A$ \unsure{faisceau de}} \inserted{Supposons $S$ affine.} Si $\mathfrak X$ est engendré par $r$ générateurs sur celui de $\mathbb W(S)$, et celui de $\mathbb W$ par $s$ gén. sur celui de $S$, alors
% archive sheet number: 67
\scanpage{69}
l'anneau de $\unsure{\underline{\mathbb W}|\underline W}^a(\unsure{\mathfrak X|\mathcal X})$ est engendré par $sr$ générateurs sur celui de $S$.

Dans le cas général \inserted{($+S$ affine)}, on a $C=\unsure{\varinjlim|\underline{\lim}} C_i$, où les $C_i$ sont des ss-$\mathbb W(A)$-algèbres \inserted{de type fini} de $C$. On a donc, pour $T$ affine
\[
% Transposition loop in the manuscript: X was written first, W(T) second.
\begin{aligned}
\mathrm{Hom}_{\underline{\mathbb W}(S)}(\underline{\mathbb W}(T),\mathfrak X)&=\mathrm{Hom}_{\mathbb W(A)\text{-alg}}(\deleted{\mathbb W(B),}C,\underline{\mathbb W}(B))\\
 &\simeq\varprojlim\mathrm{Hom}_{\underline{\mathbb W}(A)\text{-alg}}(C_i,\underline{\mathbb W}(B))\\
 &\simeq\varprojlim\mathrm{Hom}_{A\text{-alg}}(\underline{\mathbb W}^a(C_i),B)\\
 &\simeq\mathrm{Hom}_{A\text{-alg}}(\varinjlim\underline{\mathbb W}^a(C_i),B)
\end{aligned}
\]
Soit donc
\[
 \varinjlim\mathbb W^a(C_i)=\deleted{\unsure{\mathbb W^a(}}D\deleted{\unsure{,\; -}}
\]
On a obtenu un isom fonctoriel
\[
 \mathrm{Hom}_{\underline{\mathbb W}(S)}(\underline{\mathbb W}(T),\mathfrak X)\simeq\mathrm{Hom}_S(T,\operatorname{Spec}(D))
\]
d'où le résultat.

\inserted{\uline{Corollaire} Sous les conditions précédentes, si $\mathfrak X\xrightarrow{f}\unsure{\mathfrak Y|\mathcal Y}$ est un morphisme \deleted{\unsure{affine}} de $S$-préschémas affines qui est de type fini (resp. une immersion fermée) alors $\underline{\mathbb W}^a(f)$ est un morphisme de type fini \unsure{pourvu que} $\underline{\mathbb W}/S$ de type fini, \unsure{resp\textsuperscript{t.}} $\underline{\mathbb W}^a(f)$ est une immersion fermée.}
% continuation of the inserted corollary slanted in the lower left margin
\deleted{\uline{Corollaire 2} Soit $f:\mathfrak X\unsure{-|\to}\mathfrak Y$ un morphisme \uline{affine} de $\underline{\mathbb W}(S)$-préschémas, \deleted{\unsure{et}} supposons que $\underline{\mathbb W}^{\unsure{*|a}}(\mathfrak X)$ et $\underline{\mathbb W}^a(\mathfrak Y)$ existent, alors $W$ $\mathbb W^a(f):\mathbb W^a(\mathfrak X)-\mathbb W^a(\mathfrak Y)$ est affine,}
% archive sheet number: 68
\scanpage{70}
\textcircled{3} \uline{Recollement}

Aucune hypothèse sur $\unsure{\underline{\mathbb W}|\underline W}$

\uline{\deleted{Proposition} \inserted{Lemme 1}} Soit $\unsure{\mathfrak X|\mathcal X}$ sur $\underline{\mathbb W}(S)$, \deleted{$(\unsure{\mathcal U|\mathscr U|\mathfrak X}_i)_{i\in I}$ une famille d'ouverts sur $\mathfrak X$. On suppose que pour tout $T$ sur $S$, et \deleted{$\mathbb W(T)$ est \unsure{réunion}} \deleted{tout} \inserted{$\mathbb W(\underline S)$-}morphisme $\underline{\mathbb W}(T)\xrightarrow f\mathfrak X$, on a $T=\bigcup T_i$ (les $T_i$ ouverts), avec $f(\mathbb W(T_i))\subset\mathfrak X_i$. Si les $\mathbb W^a$} $\mathcal U$ un ouvert de $\mathfrak X$, On suppose que $\underline{\mathbb W}^a(\mathfrak X)=X$ existe, soit $\mathcal U'$ l'image inverse de $\mathcal U$ par
\diagram{\begin{tikzcd}[ampersand replacement=\&]
\underline{\mathbb W}(X) \arrow[r,"\unsure{\varphi|\phi}"] \& \mathfrak X \\
\mathcal U' \arrow[u,phantom,"\cup" description] \& \mathcal U \arrow[u,phantom,"\cup" description]
\end{tikzcd}}

Supposons : \deleted{qu'}\inserted{\textcircled{$\alpha$}} il existe un ouvert $U$ dans $X$ tel que pour tout $T/S$, et tout morphisme $\underline{\mathbb W}(T)\xrightarrow f\mathbb W(X)$, la condition $f(\underline{\mathbb W}(T))\subset\mathcal U'\Longleftrightarrow f(\underline{\mathbb W}(T))\subset\underline{\mathbb W}(U)$ [En particulier, on aura $\mathbb W(U)\subset\mathcal U'$]

Dans ces conditions, $U$ muni du morphisme $\mathbb W(U)\longrightarrow\mathcal U$ induit par $\underline{\mathbb W}(X)\to X$ est un $\mathbb W(\mathcal U')$.
% archive sheet number: 69
\scanpage{71}
C'est trivial. \deleted{Cas particuliers :}

\deleted{\uline{Proposition} Dans le cas \deleted{\unsure{d'un X}} I' ou II'', si $\unsure{\mathbb W|W}^a(\unsure{\mathfrak X|\mathcal X})$ existe et}

\uline{Lemme 2} Soient $\mathfrak X_i$ des ouverts de $\mathfrak X/\mathbb W(S)$, on suppose \deleted{\inserted{que $\underline{\mathbb W}^a(\mathfrak X)$ existe}} que \deleted{les $\mathfrak X_i$ dans $\mathfrak X$ satisfont les conditions du lemme 1 (en particulier} les $\mathbb W^a(\mathfrak X_i)$ existent \deleted{\unsure{et}} \deleted{\unsure{et que}} et que les $\mathfrak X_i\cap\mathfrak X_j$ dans $\mathfrak X_i$ satisfont les \deleted{\unsure{mêmes}} conditions \inserted{du lemme 1} [donc $\underline{\mathbb W}^a(\mathfrak X_i\cap\mathfrak X_j)$ existe] et est un ouvert dans $\underline{\mathbb W}^a(\mathfrak X_i)$], \inserted{et kif-kif pour les $\mathfrak X_i\cap\mathfrak X_j\cap\mathfrak X_{\unsure{h|k}}$ dans $\mathfrak X_i\cap\mathfrak X_j$} On suppose

\textcircled{$\beta$} Pour tt $T/S$, et tout morphisme $\underline{\mathbb W}(T)\xrightarrow f\mathfrak X$, on a $T=\bigcup_i T_i$ (les $T_i$ ouverts \unsure{de|dans} $T$) tel que
\[
 f(\mathbb W(T_i))\subset\mathfrak X_i
\]
Sous ces conditions, $\underline{\mathbb W}^a(\mathfrak X)$ existe, et est la réunion des $\underline{\mathbb W}^a(\mathfrak X_i)$, qui s'identifient à des ouverts de $\underline{\mathbb W}^a(\mathfrak X)$
% archive sheet number: 70
\scanpage{72}
$\unsure{\underline{\mathbb W}|\underline W}^a(\unsure{\mathfrak X|\mathcal X}_i\cap\mathfrak X_j)$ est \inserted{donc} \inserted{à la fois \unsure{un}} ouvert \deleted{à la fois} \inserted{$X_{ij}$} dans $\underline{\mathbb W}^a(\mathfrak X_i)$ et \inserted{$X_{ji}$ dans} $\mathbb W^a(\mathfrak X_j)$ : cela permet de recoller les $\mathbb W^a(\mathfrak X_i)=X_i$. Il faut vérifier pour cela la \deleted{\unsure{$X_{ij}$ et $X_{jk}$, et|$X_{ijk}=X_{ikj}$}} commutativité \unsure{de|du}.
\diagram{\begin{tikzcd}[ampersand replacement=\&,column sep=large,row sep=small]
 X_{ij\unsure{h|k}} \arrow[r,"{\unsure{\varphi|\phi}_{ji}\vert}"] \arrow[rr,bend left=25,"{\varphi_{\unsure{h|k}i}\vert}"] \& X_{jih} \arrow[r,"{\varphi_{kj}\vert}"] \& X_{kij}\\
 \parallel \& \parallel \& \parallel\\
 X_{ij}\cap X_{ih} \& X_{ji}\cap X_{jh} \& X_{ki}\cap X_{kj}
\end{tikzcd}}

Cela tient au fait que l'on aura $X_{ijh}\simeq\underline{\mathbb W}^a(\mathfrak X_i\cap\mathfrak X_j\cap\mathfrak X_k)$ et que \deleted{les} \unsure{aplication|application|applications} $\varphi_{ji}$ (p. ex) s'identifie à l'identité.

On a ainsi défini un $X=\bigcup X_i$, $X_i\simeq\underline{\mathbb W}^a(\mathfrak X_i)$. On définit
\[
 \mathbb W(X)\longrightarrow\mathfrak X
\]
en notant que $\mathbb W(X)=\bigcup_i\mathbb W(X_i)$ \deleted{\unsure{et|$\simeq$}} et qu'on a des \deleted{\unsure{hom}} morphismes
\[
 \underline{\mathbb W}(X_i)\longrightarrow\mathfrak X_i\subset\mathfrak X
\]
par définition de $X_i=\underline{\mathbb W}^a(\mathfrak X_i)$. La
% archive sheet number: 71
